% GENERATED FILE. Edit canonical lessons, then run npm run generate:books.
% canonical-source-hash: fnv1a64:16b722b52b836157
% canonical-lessons: SA-C442-suryodayah, SA-C442-udeti, SA-C442-prakasate, SA-C442-ksitijam, SA-C442-vibhati

\chapter{Sunrise}
\label{ch:sa-sunrise}
\hlchaptermodality{442}

\begin{chapteropening}
\textbf{By the end of this chapter:} \emph{I can describe the dawn and the rising sun.}

Complete the last lesson of chapter 442: I can describe the dawn and the rising sun.
\end{chapteropening}

\section[\texorpdfstring{\sk{सूर्योदयः} (sūryodayaḥ)}{सूर्योदयः (sūryodayaḥ)}]{\sk{सूर्योदयः} (sūryodayaḥ) — sunrise}

\label{lesson:SA-C442-suryodayah}



\begin{quote}
Before the new one: say the Sanskrit for hoarfrost, then the Sanskrit for veils.
\end{quote}

\subsection*{You'll want to know: \sk{सूर्योदयः}}
\textbf{\sk{सूर्योदयः}} — \emph{sūryodayaḥ} — \textquotedblleft{}sunrise\textquotedblright{}.

\begin{quote}
\sk{वयं} \sk{सूर्योदयात्} \sk{पूर्वं} \sk{प्रतिष्ठामहे।}
\end{quote}

\begin{quote}
\emph{vayaṃ sūryodayāt pūrvaṃ pratiṣṭhāmahe.}
\end{quote}

\begin{quote}
We set out before sunrise.
\end{quote}

A masculine noun.

\begin{cousinweb}
\textbf{\sk{सूर्योदयः}} (\emph{sūryodayaḥ}) is built from \emph{sūrya}, \textquotedblleft{}the sun\textquotedblright{}, and \emph{udaya}, \textquotedblleft{}rising\textquotedblright{}, from \emph{ud} and the root √i, \textquotedblleft{}to go\textquotedblright{}; MW gives the compound as sunrise (MW).
\end{cousinweb}

\begin{grammarlens}[title={What you've built}]
One more word for getting around.
\end{grammarlens}

\subsection*{Your turn}
\begin{itemize}
\raggedright
  \item \emph{Say it:} \emph{sūryodayaḥ}
  \item \emph{Say it:} \emph{sūryodayaḥ}, once more
  \item \emph{Say it:} \emph{vayaṃ sūryodayāt pūrvaṃ pratiṣṭhāmahe.} — We set out before sunrise.
  \item \emph{Recall:} say the Sanskrit for hoarfrost, then the Sanskrit for veils, then say \emph{sūryodayaḥ} again
\end{itemize}

\begin{tcolorbox}[breakable,colback=teal!4,colframe=teal!35!black,title={Before you move on}]
What does \sk{सूर्योदयः} mean? (\textquotedblleft{}Sunrise\textquotedblright{}.) Say it once more.
\end{tcolorbox}
\section[\texorpdfstring{\sk{उदेति} (udeti)}{उदेति (udeti)}]{\sk{उदेति} (udeti) — rises in the sky}

\label{lesson:SA-C442-udeti}



\begin{quote}
Before the new one: say the Sanskrit for veils, then the Sanskrit for sunrise.
\end{quote}

\subsection*{You'll want to know: \sk{उदेति}}
\textbf{\sk{उदेति}} — \emph{udeti} — \textquotedblleft{}rises in the sky\textquotedblright{}.

\begin{quote}
\sk{प्रातः} \sk{सूर्यः} \sk{उदेति।}
\end{quote}

\begin{quote}
\emph{prātaḥ sūryaḥ udeti.}
\end{quote}

\begin{quote}
In the morning the sun rises.
\end{quote}

Past (it …): \sk{सूर्यः} \sk{उदितवान्}, \emph{sūryaḥ uditavān}.

\begin{cousinweb}
\textbf{\sk{उदेति}} (\emph{udeti}) is built from \emph{ud}, \textquotedblleft{}up\textquotedblright{}, and the root √i, \textquotedblleft{}to go\textquotedblright{}, which also give \emph{udaya}, \textquotedblleft{}rising\textquotedblright{}; MW gives \emph{ud-i} as to go up and to rise, as the sun does, with the participle \emph{udita}, \textquotedblleft{}risen\textquotedblright{} (MW).
\end{cousinweb}

\begin{grammarlens}[title={What you've built}]
One more word for getting around.
\end{grammarlens}

\subsection*{Your turn}
\begin{itemize}
\raggedright
  \item \emph{Say it:} \emph{udeti}
  \item \emph{Say it:} \emph{udeti}, once more
  \item \emph{Say it:} \emph{prātaḥ sūryaḥ udeti.} — In the morning the sun rises.
  \item \emph{Recall:} say the Sanskrit for veils, then the Sanskrit for sunrise, then say \emph{udeti} again
\end{itemize}

\begin{tcolorbox}[breakable,colback=teal!4,colframe=teal!35!black,title={Before you move on}]
What does \sk{उदेति} mean? (\textquotedblleft{}Rises in the sky\textquotedblright{}.) Say it once more.
\end{tcolorbox}
\section[\texorpdfstring{\sk{प्रकाशते} (prakāśate)}{प्रकाशते (prakāśate)}]{\sk{प्रकाशते} (prakāśate) — grows bright; shines out}

\label{lesson:SA-C442-prakasate}



\begin{quote}
Before the new one: say the Sanskrit for sunrise, then the Sanskrit for rises in the sky.
\end{quote}

\subsection*{You'll want to know: \sk{प्रकाशते}}
\textbf{\sk{प्रकाशते}} — \emph{prakāśate} — \textquotedblleft{}grows bright; shines out\textquotedblright{}.

\begin{quote}
\sk{सूर्योदये} \sk{आकाशः} \sk{प्रकाशते।}
\end{quote}

\begin{quote}
\emph{sūryodaye ākāśaḥ prakāśate.}
\end{quote}

\begin{quote}
At sunrise the sky grows bright.
\end{quote}

Past (it …): \sk{आकाशः} \sk{प्रकाशितवान्}, \emph{ākāśaḥ prakāśitavān}.

\begin{cousinweb}
\textbf{\sk{प्रकाशते}} (\emph{prakāśate}) is built from \emph{pra}, \textquotedblleft{}forth\textquotedblright{}, and the root √kāś, \textquotedblleft{}to shine\textquotedblright{}, which also gives \emph{prakāśa}, \textquotedblleft{}light\textquotedblright{}; MW gives \emph{pra-kāś} as to become visible, to shine, with the participle \emph{prakāśita} (MW).
\end{cousinweb}

\begin{grammarlens}[title={What you've built}]
One more word for getting around.
\end{grammarlens}

\subsection*{Your turn}
\begin{itemize}
\raggedright
  \item \emph{Say it:} \emph{prakāśate}
  \item \emph{Say it:} \emph{prakāśate}, once more
  \item \emph{Say it:} \emph{sūryodaye ākāśaḥ prakāśate.} — At sunrise the sky grows bright.
  \item \emph{Recall:} say the Sanskrit for sunrise, then the Sanskrit for rises in the sky, then say \emph{prakāśate} again
\end{itemize}

\begin{tcolorbox}[breakable,colback=teal!4,colframe=teal!35!black,title={Before you move on}]
What does \sk{प्रकाशते} mean? (\textquotedblleft{}Grows bright; shines out\textquotedblright{}.) Say it once more.
\end{tcolorbox}
\section[\texorpdfstring{\sk{क्षितिजम्} (kṣitijam)}{क्षितिजम् (kṣitijam)}]{\sk{क्षितिजम्} (kṣitijam) — the horizon}

\label{lesson:SA-C442-ksitijam}



\begin{quote}
Before the new one: say the Sanskrit for rises in the sky, then the Sanskrit for grows bright.
\end{quote}

\subsection*{You'll want to know: \sk{क्षितिजम्}}
\textbf{\sk{क्षितिजम्}} — \emph{kṣitijam} — \textquotedblleft{}the horizon\textquotedblright{}.

\begin{quote}
\sk{प्रातः} \sk{क्षितिजम्} \sk{अरुणं} \sk{भवति।}
\end{quote}

\begin{quote}
\emph{prātaḥ kṣitijam aruṇaṃ bhavati.}
\end{quote}

\begin{quote}
In the morning the horizon turns red.
\end{quote}

A neuter noun.

\begin{cousinweb}
\textbf{\sk{क्षितिजम्}} (\emph{kṣitijam}) is built from \emph{kṣiti}, \textquotedblleft{}the earth\textquotedblright{}, and \emph{ja}, \textquotedblleft{}born\textquotedblright{}: MW gives \emph{kṣitija} as earth-born and, as a neuter, the horizon (MW).
\end{cousinweb}

\begin{grammarlens}[title={What you've built}]
One more word for getting around.
\end{grammarlens}

\subsection*{Your turn}
\begin{itemize}
\raggedright
  \item \emph{Say it:} \emph{kṣitijam}
  \item \emph{Say it:} \emph{kṣitijam}, once more
  \item \emph{Say it:} \emph{prātaḥ kṣitijam aruṇaṃ bhavati.} — In the morning the horizon turns red.
  \item \emph{Recall:} say the Sanskrit for rises in the sky, then the Sanskrit for grows bright, then say \emph{kṣitijam} again
\end{itemize}

\begin{tcolorbox}[breakable,colback=teal!4,colframe=teal!35!black,title={Before you move on}]
What does \sk{क्षितिजम्} mean? (\textquotedblleft{}The horizon\textquotedblright{}.) Say it once more.
\end{tcolorbox}
\section[\texorpdfstring{\sk{विभाति} (vibhāti)}{विभाति (vibhāti)}]{\sk{विभाति} (vibhāti) — dawns; shines forth}

\label{lesson:SA-C442-vibhati}



\begin{quote}
Before the new one: say the Sanskrit for grows bright, then the Sanskrit for the horizon.
\end{quote}

\subsection*{You'll want to know: \sk{विभाति}}
\textbf{\sk{विभाति}} — \emph{vibhāti} — \textquotedblleft{}dawns; shines forth\textquotedblright{}.

\begin{quote}
\sk{शनैः} \sk{रात्रिः} \sk{विभाति।}
\end{quote}

\begin{quote}
\emph{śanaiḥ rātriḥ vibhāti.}
\end{quote}

\begin{quote}
Slowly the night turns to dawn.
\end{quote}

Past (it …): \sk{रात्रिः} \sk{विभातवती}, \emph{rātriḥ vibhātavatī}.

\begin{cousinweb}
\textbf{\sk{विभाति}} (\emph{vibhāti}) is built from \emph{vi} and the root √bhā, \textquotedblleft{}to shine\textquotedblright{}; MW gives \emph{vi-bhā} as to shine forth, and its participle \emph{vibhāta} for a night that has grown light, the morning that has dawned (MW).
\end{cousinweb}

\begin{grammarlens}[title={What you've built}]
One more word for getting around.
\end{grammarlens}

\subsection*{Your turn}
\begin{itemize}
\raggedright
  \item \emph{Say it:} \emph{vibhāti}
  \item \emph{Say it:} \emph{vibhāti}, once more
  \item \emph{Say it:} \emph{śanaiḥ rātriḥ vibhāti.} — Slowly the night turns to dawn.
  \item \emph{Recall:} say the Sanskrit for grows bright, then the Sanskrit for the horizon, then say \emph{vibhāti} again
\end{itemize}

\begin{tcolorbox}[breakable,colback=teal!4,colframe=teal!35!black,title={Before you move on}]
What does \sk{विभाति} mean? (\textquotedblleft{}Dawns; shines forth\textquotedblright{}.) Say it once more.
\end{tcolorbox}
